\section{Exact normalized PPR geometry}
\label{sec:normalized-geometry}

The first information-level question can be answered without an algorithm.
Define the degree-normalized exact PPR potential and its superlevel set by
\begin{equation}
 y^0:=D^{-1/2}x^0=D^{-1}\pi,
 \qquad
 U_\tau^0:=\{i\in V:y_i^0>\tau\}.
 \label{eq:normalized-potential}
\end{equation}
Multiplying $Qx^0=b$ by $D^{1/2}$ and substituting
$x^0=D^{1/2}y^0$ gives the degree-form system
\begin{equation}
 H_\alpha y^0=\alpha e_v,
 \qquad
 H_\alpha:=D^{1/2}QD^{1/2}
 =\alpha D+\frac{1-\alpha}{2}(D-A).
 \label{eq:killed-system}
\end{equation}

\begin{lemma}[Strict ascent and seed-connected superlevels; \ProvedStatus]
\label{lem:strict-ascent}
For every $i\neq v$,
\begin{equation}
 y_i^0
 =\frac{1-\alpha}{1+\alpha}
   \frac1{d_i}\sum_{j\in\mathcal N(i)}y_j^0.
 \label{eq:discounted-average}
\end{equation}
If $y_i^0>0$, then $i$ has a neighbor $j$ with $y_j^0>y_i^0$.
Consequently $y_v^0=\max_i y_i^0$, every nonempty $U_\tau^0$ contains $v$,
and every vertex of $U_\tau^0$ has a path to $v$ lying entirely in
$U_\tau^0$.
\end{lemma}

\begin{proof}
The nonsource row of \cref{eq:killed-system} is
\[
 \frac{1+\alpha}{2}d_i y_i^0
 =\frac{1-\alpha}{2}\sum_{j\in\mathcal N(i)}y_j^0,
\]
which is \cref{eq:discounted-average}.  Its prefactor is strictly below one,
so a positive $y_i^0$ is strictly smaller than at least one neighboring
value.  Repeating this strict ascent cannot cycle in a finite graph and can
terminate only at the sourced row $v$.  Starting from a value above $\tau$
keeps every vertex of the ascent path above $\tau$.
\end{proof}

\begin{theorem}[Exact information core; \ProvedStatus]
\label{thm:information-core}
For every $\tau>0$,
\begin{equation}
 \vol(U_\tau^0)\leq\frac1\tau.
 \label{eq:information-core-volume}
\end{equation}
Moreover, the omniscient truncation
\begin{equation}
 \widehat x_i:=
 \begin{cases}
 x_i^0,&i\in U_\tau^0,\\
 0,&i\notin U_\tau^0
 \end{cases}
 \label{eq:omniscient-truncation}
\end{equation}
satisfies
$\norm{D^{-1/2}(\widehat x-x^0)}_\infty\leq\tau$ and has support degree
volume at most $1/\tau$.
\end{theorem}

\begin{proof}
Left-multiplying \cref{eq:killed-system} by $\one^\top$ and using
$\one^\top(D-A)=0$ yields
\begin{equation}
 \sum_i d_i y_i^0=1.
 \label{eq:normalized-mass}
\end{equation}
Therefore
$1\geq\sum_{i\in U_\tau^0}d_i y_i^0
>\tau\vol(U_\tau^0)$.
Outside $U_\tau^0$, the normalized truncation error equals $y_i^0\leq\tau$;
inside it is zero.
\end{proof}

The theorem is representational, not algorithmic: membership in
$U_\tau^0$ and its exact values are initially unknown.  It strengthens the
existential interpretation of the exact Dirichlet-core theorem in the
companion \texttt{volume\_gated\_acceleration} note.  Strict ascent also gives
the useful distance consequence
\begin{equation}
 i\in U_\tau^0
 \quad\Longrightarrow\quad
 \operatorname{dist}(v,i)+1
 \leq |U_\tau^0|
 \leq\vol(U_\tau^0)
 \leq\frac1\tau.
 \label{eq:significant-distance}
\end{equation}

\begin{theorem}[Tight explicit-output complexity; \ProvedStatus]
\label{thm:output-tight}
Fix $0<\alpha\leq1/2$ and
$0<\epsppr\leq(1-\alpha)/8$.  In the word-output model, the worst-case
number of listed coordinates required by \cref{eq:note-ppr-error} is
\begin{equation}
 \Theta(1/\epsppr).
 \label{eq:output-theta}
\end{equation}
The conclusion concerns output words; bit precision introduces the usual
additional representation factors.
\end{theorem}

\begin{proof}
The upper bound follows from \cref{thm:information-core}, because every
nonisolated vertex has degree at least one.  For the lower bound, take a star
with seed $v$ at its center and
\[
 M:=\left\lfloor\frac{1-\alpha}{4\epsppr}\right\rfloor
\]
leaves.  Direct substitution into \cref{eq:killed-system} gives
\begin{equation}
 y_v^0=\frac{1+\alpha}{2M},
 \qquad
 y_i^0=\frac{1-\alpha}{2M}
 \quad\text{for every leaf }i.
 \label{eq:star-normalized-values}
\end{equation}
Thus every leaf has $y_i^0\geq2\epsppr$.  An output satisfying
\cref{eq:note-ppr-error} must list a nonzero approximation at every leaf, and
$M=\Omega(1/\epsppr)$ in the stated regime.
\end{proof}

The best currently documented general upper bounds should not be confused
with this existential result.  Wei and Yang's growing-active-set method has
an ACL upper bound
$\widetilde O(\min\{1/\varepsilon^2,1/(\alpha\varepsilon)\})$ in its native,
nonlazy convention.  Their Section~6, PDF pages 14--15, explicitly observes
that this rules out uniform
$\Omega(1/(\alpha\varepsilon))$ and
$\Omega(1/(\sqrt\alpha\varepsilon))$ lower bounds across all parameter
regimes, and identifies reuse toward $\widetilde O(1/\varepsilon)$ as open
\citep{wei2026simple}.  After the teleportation mapping recorded in the
companion \texttt{incremental\_active\_set\_sdd} note, this is fully
consistent with \cref{thm:output-tight}; it does not yet prove an
$\widetilde O(1/\epsppr)$ adjacency-query algorithm on arbitrary graphs.
